% Preámbulo compartido por presentacion.tex (diapositivas) y guion.tex (notas de cada orador).
\documentclass[aspectratio=1610]{beamer}

% ---------- Tema estilo "sidebar" (Berkeley) en verde oscuro, como en la Parte 1 ----------
\usetheme{Berkeley}
\usecolortheme{default}

\usepackage[spanish]{babel}
\usepackage[utf8]{inputenc}
\usepackage[T1]{fontenc}
\usepackage{booktabs}
\usepackage{graphicx}
\usepackage{amsmath}
\usepackage{array}
\usepackage{multirow}

\usepackage{tikz}
\usetikzlibrary{shadows.blur, arrows.meta}
\usepackage{pgfplots}
\pgfplotsset{compat=1.18}

\graphicspath{{../Parte_1/Figuras/}{../Parte_2/Figuras/}{../Parte_2/Planos/}}

% Paleta verde oscuro
\definecolor{unc verde}{RGB}{27,58,48}
\definecolor{unc verdeclaro}{RGB}{60,110,90}
\definecolor{grupo este}{RGB}{27,58,48}
\definecolor{grupo oeste}{RGB}{120,60,30}

\setbeamercolor{palette primary}{bg=unc verde,fg=white}
\setbeamercolor{palette secondary}{bg=unc verde,fg=white}
\setbeamercolor{palette tertiary}{bg=unc verde,fg=white}
\setbeamercolor{palette quaternary}{bg=unc verde,fg=white}
\setbeamercolor{structure}{fg=unc verde}
\setbeamercolor{title}{bg=unc verde,fg=white}
\setbeamercolor{frametitle}{bg=unc verde,fg=white}
\setbeamercolor{sidebar}{bg=unc verde}
\setbeamercolor{section in sidebar}{fg=white}
\setbeamercolor{section in sidebar shaded}{fg=unc verdeclaro}
\setbeamercolor{subsection in sidebar}{fg=white}
\setbeamercolor{subsection in sidebar shaded}{fg=unc verdeclaro}
\setbeamercolor{title in sidebar}{fg=white}
\setbeamercolor{author in sidebar}{fg=white}
\setbeamerfont{section in sidebar}{series=\bfseries, size=\fontsize{5.5}{6.5}\selectfont}

\title[TP4: Fibra óptica]{Actividad 4: Fibra óptica}
\subtitle{Parte 1: enlace P2P -- Parte 2: red FTTH GPON}
\author[Arteaga, Baccino, Mendes Rosa, Monja]{\texorpdfstring{Arteaga Barrera, Cristian Eduardo \\ Baccino, Luca \\ Mendes Rosa, Agustín \\ Monja, Ernesto Joaquín}{Arteaga Barrera, Cristian Eduardo; Baccino, Luca; Mendes Rosa, Agustín; Monja, Ernesto Joaquín}}
\institute[FCEFyN -- UNC]{Facultad de Ciencias Exactas, Físicas y Naturales \\ Universidad Nacional de Córdoba}
\date{Córdoba, 2026}

% ---------- Orador de cada diapositiva ----------
% \orador{Grupo}{Nombre} fija quién expone desde esa diapositiva en adelante; se muestra en el pie.
\newcommand{\oradorgrupo}{}
\newcommand{\oradornombre}{}
\newcommand{\oradorcolor}{grupo este}
\newcommand{\orador}[2]{%
  \gdef\oradorgrupo{#1}\gdef\oradornombre{#2}%
  \ifnum\pdfstrcmp{#1}{Grupo Oeste}=0 \gdef\oradorcolor{grupo oeste}\else\gdef\oradorcolor{grupo este}\fi}

\setbeamertemplate{navigation symbols}{}
% El orador no se muestra en las diapositivas: figura solo en el guion (guion.tex).
\setbeamertemplate{footline}{%
  \hfill{\scriptsize\color{unc verde}\insertframenumber}\hspace{2mm}\vspace{1.5mm}}

% Notas del guion: tipografía legible
\setbeamerfont{note page}{size=\small}
